\documentclass[11pt]{article}

\usepackage[margin=0.85in]{geometry}
\usepackage[T1]{fontenc}
\usepackage[utf8]{inputenc}
\usepackage{amsmath,amssymb}
\usepackage{xcolor}
\usepackage{tcolorbox}
\usepackage{enumitem}
\usepackage{hyperref}

\definecolor{auditred}{RGB}{185,38,38}
\definecolor{auditblue}{RGB}{34,80,145}
\definecolor{auditgray}{RGB}{245,246,248}

\setlength{\parindent}{0pt}
\setlength{\parskip}{0.55em}
\setlist[itemize]{leftmargin=1.4em,itemsep=0.2em,topsep=0.2em}

\newtcolorbox{issuebox}[1]{
  colback=auditgray,
  colframe=#1,
  boxrule=0.8pt,
  arc=1.5pt,
  left=7pt,
  right=7pt,
  top=6pt,
  bottom=6pt
}

\begin{document}

\begin{center}
{\Large\bfseries Minimal sign issue in the shortcut-to-absorbing-site calculation}\\[0.35em]
{\small Note on the first hard error in Eqs. (30)--(32)}
\end{center}

\begin{issuebox}{auditblue}
\textbf{Where I think the issue starts.}
The step from the probability-conserving shortcut formula to the absorbing
case is made by replacing the ring propagator $\widetilde P$ with the absorbing
propagator $\widetilde W$.  This replacement is not valid once the shortcut
endpoint is also the absorbing site.
\end{issuebox}

In the manuscript, the absorbing site is then set equal to the shortcut
endpoint:
\[
m=v,\qquad \widetilde W_v(n,z)=0.
\]
At this point, the shortcut $u\to v$ is not an additional transient propagation
channel.  Since $v$ is absorbing, a walker that takes the shortcut has left the
transient state space.

Let
\[
\lambda=\beta(1-q),
\]
and let $P_0$ denote the transition matrix on the non-absorbing states for the
ring with the absorbing target $v$ already removed.  In transient space, the
shortcut to $v$ is therefore a killing term at $u$:
\[
\boxed{P=P_0-\lambda |u\rangle\langle u|.}
\]
Consequently
\[
I-zP=I-zP_0+z\lambda |u\rangle\langle u|.
\]
Applying Sherman--Morrison gives
\[
(I-zP)^{-1}
=(I-zP_0)^{-1}
-\frac{
z\lambda (I-zP_0)^{-1}|u\rangle\langle u|(I-zP_0)^{-1}
}{
1+z\lambda\langle u|(I-zP_0)^{-1}|u\rangle
}.
\]

\begin{issuebox}{auditred}
\textbf{Corrected form of Eq. (31).}
\[
\boxed{
\widetilde S_{n_0}(n,z)
=\widetilde W_{n_0}(n,z)
-\frac{
z\beta(1-q)\widetilde W_{n_0}(u,z)\widetilde W_u(n,z)
}{
1+z\beta(1-q)\widetilde W_u(u,z)
}.}
\]
The correction term is negative and the denominator has a plus sign.
\end{issuebox}

The sign is forced by the meaning of the absorbing boundary: adding a direct
shortcut into the absorbing target removes paths from the transient propagator.
Equivalently, to first order in $\lambda$ the transient occupation should
decrease:
\[
\widetilde S
=\widetilde W
-z\lambda\,\widetilde W_{n_0}(u,z)\widetilde W_u(n,z)
+O(\lambda^2).
\]
The formula with a positive correction and denominator
$1-z\lambda\widetilde W_u(u,z)$ would instead increase transient occupation,
which corresponds to adding a propagation channel inside the transient state
space rather than adding absorption at $v$.

The first-passage generating function then inherits the corrected sign:
\[
\boxed{
\widetilde F_{n_0}(v,z)
=\widetilde{\mathcal F}_{n_0}(v,z)
+\frac{
z\beta(1-q)\widetilde W_{n_0}(u,z)
\left[1-\widetilde{\mathcal F}_u(v,z)\right]
}{
1+z\beta(1-q)\widetilde W_u(u,z)
}.}
\]

Thus the later auxiliary denominator should be
\[
\frac{\beta(1-q)}{q}\,
\frac{1}{1+z\beta(1-q)\widetilde W_u(u,z)},
\]
not the corresponding expression with $1-z\beta(1-q)\widetilde W_u(u,z)$.
The later pole and threshold conclusions should therefore be recomputed from
this corrected resolvent.

\end{document}
